\section{Conclusion}

This study examined the automatic prioritization of IT support tickets using
NLP and machine learning together with the SHAP-based explainability of the
resulting model predictions. Seven classifiers were compared under a common
group-based evaluation design and assessed with respect to classification
performance and selected explanation properties.

RQ1 shows that the investigated classifiers differed in their prioritization
performance, although no general superiority of a particular model family was
observed. On FS4, which was selected as the common feature set based on
cross-validation, LightGBM achieved the highest hold-out performance and was
the only alternative that significantly outperformed Logistic Regression after
Holm correction. H1 is therefore supported. The final model selected overall
was LightGBM with FS5, which achieved a Macro-F1 of 0.5597 on the separate
hold-out dataset. In addition to the ticket text, the queue provided relevant
contextual information, whereas the additional contribution of sentiment
features was small.

RQ2 also revealed clear model-dependent differences. Logistic Regression
showed significantly higher SHAP-based explanation compactness than Complement
Naive Bayes and Random Forest, but not than all alternative classifiers. H2 is
therefore supported according to the prespecified criterion. At the same time,
LightGBM combined comparatively high classification performance with a high
concentration of SHAP attributions. The results consequently provide no
evidence of a consistent trade-off between predictive performance and
explanation compactness.

The interpretation of these findings is mainly limited by the synthetic data
basis, the absence of independent validation of the priority labels, and the
evaluation within a single dataset. Further limitations arise from the
organization-specific role of the queue and from methodological constraints
of the SHAP-based model comparisons. The findings should therefore be
understood as a controlled comparison of the investigated approaches rather
than as evidence of readiness for operational deployment.

Future research should replicate the analysis on real-world ITSM data,
preferably across multiple organizations, and use temporal test sets to examine
concept drift. Ordinal and cost-sensitive learning approaches may offer useful
extensions, as may multilingual transformer-based models such as
XLM-R \cite{conneau2020xlmr}. Explainability research should also examine whether the
identified SHAP explanations are perceived by ITSM experts as understandable
and relevant to decision making. A subsequent field study could then assess
whether model-supported prioritization in a human-in-the-loop setting leads to
measurable improvements in operational service metrics.